\section{为什么本书需要一个遵守研究纪律的案例库}

如果没有案例，这套框架就会过于接近一篇有说服力的议论文。\nobreakspace 如果没有证据纪律，案例又容易退化为语气笃定的讲故事。因而，本章刻意占据一个中间层。它\emph{并不}声称，日常的学习失败、锻炼回避或睡前拖延，已经都被一一对应地用神经科学精确测量出来。它所主张的内容更窄，但也更有用：
\begin{enumerate}[nosep]
  \item 框架中的各个变量提供了一种稳定的分析语言；
  \item 来自经验层、涌现层和高级动力学章节的研究主干，为机制类比提供了有纪律的支撑；
  \item 当应用层干预被明确系于状态转变结构，而不是模糊的动机说辞时，它就会更强。
\end{enumerate}

因此，本章对每个案例都采用三车道阅读规则：
\begin{itemize}[nosep]
  \item \textbf{框架车道：} 哪些变量占主导？
  \item \textbf{机制车道：} 哪些有研究支撑的状态、路径依赖、亚稳态或吸引子观念与之相关？
  \item \textbf{证据车道：} 哪些判断属于强证据，哪些属于中等证据，哪些一旦推得过远就会滑向推测？
\end{itemize}